\chapter{Branching and charge tables}
\label{app:branching}

\section{One family of the chiral block}
\label{app:branching:family}

The block $V_*=\operatorname{im}P_*$ carries $(\mathbf4,\mathbf2,\mathbf1)_L\oplus(\mathbf4,\mathbf1,\mathbf2)_R$ of
$\su4\oplus\su2_L\oplus\su2_R$ (\cref{thm:PS}; proved, checks in \cref{app:manifest}). Its decomposition under
$\su3\oplus\su2_L\oplus\uu1_Y$, with $\mathbf4=\mathbf3_{1/3}\oplus\mathbf1_{-1}$ under $\su3\oplus\uu1_{B-L}$ and
$Q=T^L_3+T^R_3+\tfrac12(B-L)$, is the standard one \cite{PatiSalam}; it is inherited, and the charge operators of the block reproduce it
(\cref{ch:ch08_photon}).

\begin{table}[ht]
\centering\small
\begin{tabular}{@{}lccccc@{}}
\toprule
Pati--Salam component & states & $\su3$ & $T^L_3$ or $T^R_3$ & $B-L$ & $Q$\\
\midrule
$(\mathbf4,\mathbf2,\mathbf1)_L$ & $u_L$, $d_L$ & $\mathbf3$ & $\pm\tfrac12$ & $\tfrac13$ & $\tfrac23$, $-\tfrac13$\\
 & $\nu_L$, $e_L$ & $\mathbf1$ & $\pm\tfrac12$ & $-1$ & $0$, $-1$\\
$(\mathbf4,\mathbf1,\mathbf2)_R$ & $u_R$, $d_R$ & $\mathbf3$ & $\pm\tfrac12$ & $\tfrac13$ & $\tfrac23$, $-\tfrac13$\\
 & $\nu_R$, $e_R$ & $\mathbf1$ & $\pm\tfrac12$ & $-1$ & $0$, $-1$\\
\bottomrule
\end{tabular}
\caption{One family in the chiral block: $16$ complex, $64$ real dimensions with the Lorentz spinor components.}
\label{tab:app:family}
\end{table}

\section{The space of mass kernels}
\label{app:branching:kernels}

With the complex structure $\widetilde{\mathcal C}$ of \cref{thm:kernelspace}, the $272$-dimensional real space $\mathcal M$ of Lorentz-invariant mass
kernels is a complex space $V$ of dimension $136$, and $\mathcal M\otimes\C=V\oplus\bar V$. \Cref{tab:app:kernels} lists its
components (proved: \cref{thm:kernelspace}; the invariant counts built on them are checked exactly by z4, the conjugation of the tens by the cubic
trace in a11-1). Replacing $\widetilde{\mathcal C}$ by $-\widetilde{\mathcal C}$ conjugates every label. Fermion number acts with charge $2$ on the Majorana-type components and $0$ on the Dirac-type
ones (\cref{prop:fermionnumber}).

\begin{table}[ht]
\centering\small
\begin{tabular}{@{}lccc@{}}
\toprule
component of $V$ & complex dimension & type & fermion number\\
\midrule
$(\bar{\mathbf{10}},\mathbf3,\mathbf1)$ & $30$ & Majorana & $2$\\
$(\mathbf{10},\mathbf1,\mathbf3)$ & $30$ & Majorana & $2$\\
$(\mathbf{15},\mathbf2,\mathbf2)$ & $60$ & Dirac & $0$\\
$2\,(\mathbf6,\mathbf1,\mathbf1)$ & $12$ & Majorana & $2$\\
$(\mathbf1,\mathbf2,\mathbf2)$ & $4$ & Dirac & $0$\\
\midrule
total & $136$ & & \\
\bottomrule
\end{tabular}
\caption{The mass kernels of the chiral block. In the language of $\so{10}$ the content is that of $\mathbf{10}\oplus\overline{\mathbf{126}}$, the dual of
$\mathrm{Sym}^2(\mathbf{16})=\mathbf{10}\oplus\mathbf{126}$: $\mathbf{10}=(\mathbf6,\mathbf1,\mathbf1)\oplus(\mathbf1,\mathbf2,\mathbf2)$ and
$\overline{\mathbf{126}}=(\mathbf6,\mathbf1,\mathbf1)\oplus(\overline{\mathbf{10}},\mathbf3,\mathbf1)\oplus(\mathbf{10},\mathbf1,\mathbf3)\oplus(\mathbf{15},\mathbf2,\mathbf2)$
(inherited), with $\mathbf{16}=(\mathbf4,\mathbf2,\mathbf1)\oplus(\bar{\mathbf4},\mathbf1,\mathbf2)$ and the complex structure $\widetilde{\mathcal C}$ of \cref{thm:kernelspace}.}
\label{tab:app:kernels}
\end{table}
